\section{Exact Gamma moments and their endpoint terms}
\label{mom:section}

The coefficient theorem has an immediate analytic use, but it is
essential to keep the complete normal form and its primitive.
We give the inherited scalar bridge with its normalization, then an
explicit ninth-weight identity containing the first cube term.
The bridge itself is already proved in \cite[Section 4]{PTI}.

\subsection{Height-one Laurent coefficients as ordinary integrals}

Define rational zeta polynomials $c_j$ and height-one Laurent
coefficients $\lambda_d$ by
\begin{equation}\label{mom:clambda}
 \frac1{\Gamma(1+u)}=\sum_{j\ge0}c_ju^j,\qquad
 \lambda_d=[u^1]\zeta(1+u,\{1\}^{d-1}),\qquad d\ge1.
\end{equation}
The coefficient on the right is a Laurent coefficient: for $d>1$
the pole at zero has order $d$. In particular
$c_0=1$, $c_1=\gamma$, and $\lambda_1=-\gamma_1$.
For $v>0$ put
\[
 L(v)=-\log(1-e^{-v}),\qquad
 V_d(v)=L(v)^d-\boldsymbol1_{(0,1)}(v)(-\log v)^d.
\]
Then the normalization can be stated without an unevaluated finite
part:
\begin{equation}\label{mom:lambdaintegral}
 \boxed{\displaystyle
 \lambda_d=c_{d+1}+\frac1{d!}
                    \int_0^\infty\frac{V_d(v)}v\dd v.}
\end{equation}
The integral is ordinary and absolutely convergent. At zero its
integrand is $O((1+|\log v|)^{d-1})$; at infinity it decays
exponentially up to the factor $1/v$.

Here is a direct derivation. With $H_{n-1}(\{1\}^{d-1})$ denoting
the strict unit-index multiple harmonic sum, its generating function
is
\[
 \sum_{n\ge1}H_{n-1}(\{1\}^{d-1})e^{-nv}
 =-\frac1{d!}\frac{d}{dv}L(v)^d.
\]
This follows by differentiating the power series for $L(v)^d$,
or inductively from strict summation. Mellin integration and
integration by parts, initially for $\re u>0$, give
\begin{align}\label{mom:heightmellin}
 \zeta(1+u,\{1\}^{d-1})
 &=\frac1{d!\Gamma(u)}\int_0^\infty v^{u-1}L(v)^d\dd v\notag\\
 &=\frac1{\Gamma(1+u)u^d}
   +\frac{u}{d!\Gamma(1+u)}
          \int_0^\infty v^{u-1}V_d(v)\dd v.
\end{align}
The last integral is holomorphic for $u$ in a neighborhood of zero.
Extracting $[u^1]$ proves \eqref{mom:lambdaintegral} and fixes its sign.

\subsection{The scalar inverse-Gamma evaluation}

Let $\Acal_m$ be the differential polynomials of
\eqref{var:hierarchy}. For $m\ge2$ set
\[
 \mathcal L_m=\FP\int_0^1\Acal_m(x)\dd x,
 \qquad K_{2r}=\frac{(-1)^r\pi^{2r}}{2r+1},\quad K_{2r+1}=0.
\]
The inherited all-order scalar formula is
\begin{equation}\label{mom:scalar}
 \boxed{\displaystyle
 \mathcal L_m=K_m-m!\sum_{d=1}^{m-1}
                (-1)^{m-1-d}c_{m-1-d}\lambda_d.}
\end{equation}
We recall why no coordinate constant is suppressed. Put
\[
 \mathcal C_m(x)=m![t^m]
                  \frac{\Gamma(1-t)\Gamma(x)}{\Gamma(x-t)}.
\]
For $m\ge2$, $\mathcal C_m$ extends continuously to $[0,1]$.
Indeed the residue of the displayed quotient at $x=0$ is exactly
$-t$, with no higher powers of $t$. Euler's beta integral applied to
$F(x,t)=\Gamma(1+t)\Gamma(x)/\Gamma(x+t)$ yields
\begin{align}\label{mom:Phi}
 \int_0^1\{F(x,t)-1+t(\psi(x)+\gamma)\}\dd x
 =t\int_0^\infty
 \frac{(1-e^{-v})^t-1
       -\boldsymbol1_{(0,1)}(v)(v^t-1)}v\dd v.
\end{align}
For clarity, the subtraction follows by integrating in $x$ first
when $\re t>0$. The resulting expression is
\[
 t\int_0^\infty
 \left\{\frac{(1-e^{-v})^t-1}{v}
             +\frac{e^{-v}}{1-e^{-v}}\right\}\dd v-1.
\]
At lower cutoff $\delta$, the second integral is
$-\log(1-e^{-\delta})$, while
$\int_\delta^1(v^t-1)\dd v/v=(1-\delta^t)/t+\log\delta$.
Their limits prove \eqref{mom:Phi}. Its subtracted integrand also
justifies holomorphic continuation and coefficient extraction near
$t=0$.

Replacing $t$ by $-t$ and using \eqref{mom:lambdaintegral} now gives
\[
 \int_0^1\mathcal C_m(x)\dd x
 =-m\int_0^\infty\frac{V_{m-1}(v)}v\dd v
 =m!(c_m-\lambda_{m-1}),\qquad m\ge2.
\]
The constant and linear coefficientwise finite parts of the normalized
quotient are $1$ and $\gamma$, since
$\mathcal C_1=\psi+\gamma$ and $\FP\int_0^1\psi(x)\dd x=0$.
Thus its complete coefficientwise integral is
\[
 \frac1{\Gamma(1+t)}-\sum_{d\ge1}\lambda_dt^{d+1}.
\]
Multiplication by $1/\Gamma(1-t)$ and the reflection identity
$1/(\Gamma(1+t)\Gamma(1-t))=\sin(\pi t)/(\pi t)$ prove
\eqref{mom:scalar}. All uses of an infinite generating series here
are coefficientwise identities justified by the displayed local
subtractions.

\subsection{The complete differential identity and a cube-sector moment}

Write $P_m=\NN(\Acal_m)$ and $H_m=\HH(\Acal_m)$. With
\begin{equation}\label{mom:Delta}
 \Delta H=H(\psi(1),\psi'(1),\ldots)
        -[x^0]H(\psi(x),\psi'(x),\ldots),
\end{equation}
the exact transfer is
\begin{equation}\label{mom:transfer}
 \boxed{\displaystyle
 \FP\int_0^1P_m(\psi,\psi',\ldots)\dd x
 =K_m-m!\sum_{d=1}^{m-1}(-1)^{m-1-d}c_{m-1-d}\lambda_d
       -\Delta H_m.}
\end{equation}
This follows from $\Acal_m=P_m+DH_m$ and
\eqref{scope:boundary}. The endpoint is finite algebra: the standard
local series \cite[Section 5.7]{DLMF} is
\[
 \psi(x)=-x^{-1}-\gamma+
             \sum_{j\ge1}(-1)^{j+1}\zeta(j+1)x^j.
\]
For a monomial $M=\prod u_r^{a_r}$ of weight
$w=\sum(r+1)a_r$, multiply its substitutions by $x^w$ and take
the coefficient of $x^w$. Only the first $w+1$ coefficients of each
normalized series $x^{r+1}\psi^{(r)}(x)$ can contribute. This gives
a terminating algorithm in $\Q[\gamma,\zeta(2),\ldots,\zeta(w)]$.

For example, the full ninth-weight normal polynomial is
\begin{equation}\label{mom:P9}
 P_9=u_0^9-126u_0^5u_1^2+84u_0^3u_2^2
       -63u_0u_1^4-9u_0u_3^2+10u_2^3.
\end{equation}
The coefficient $10$ is the first cube term from
Corollary~\ref{var:squarecube}. For a fully explicit endpoint write
$Z_j=\zeta(j)$ and define
\begin{align}\label{mom:B9}
 B_9={}&-252\gamma^6Z_2+1260\gamma^5Z_3
          +\gamma^4(3150Z_2^2-5670Z_4)\notag\\
 &+\gamma^3(-17640Z_2Z_3+20328Z_5)\notag\\
 &+\gamma^2(-7560Z_2^3+40572Z_2Z_4+16884Z_3^2-57456Z_6)\notag\\
 &+\gamma(28476Z_2^2Z_3-70560Z_2Z_5-54432Z_3Z_4+106416Z_7)\notag\\
 &+1932Z_2^4-22050Z_2^2Z_4-18648Z_2Z_3^2+63036Z_2Z_6\notag\\
 &+47376Z_3Z_5+21942Z_4^2-99624Z_8.
\end{align}
Then $B_9=\Delta H_9$, and the concrete identity is
\begin{align}\label{mom:ninthidentity}
 \FP\int_0^1\bigl\{&\psi^9-126\psi^5(\psi')^2
   +84\psi^3(\psi'')^2-63\psi(\psi')^4
   -9\psi(\psi''')^2+10(\psi'')^3\bigr\}\dd x\notag\\
 &=-9!\sum_{d=1}^8(-1)^{8-d}c_{8-d}\lambda_d-B_9.
\end{align}
All functions inside the integral have argument $x$.
The formula specifies an evaluated combination of mixed moments;
it is not a claim that each term separately reduces to ordinary zeta
values. If desired, every $\lambda_d$ on its right can be replaced
by the convergent integral \eqref{mom:lambdaintegral}.

For comparison, the next full normal polynomial is
\begin{align}\label{mom:P10}
 P_{10}={}&u_0^{10}-210u_0^6u_1^2+210u_0^4u_2^2
       -315u_0^2u_1^4-45u_0^2u_3^2\notag\\
       &+100u_0u_2^3+90u_1^2u_2^2+u_4^2.
\end{align}
The companion exact data file gives $P_m,H_m$ through weight $14$,
both endpoint polynomials at weights $9,10$, and the polynomial
numerators $p_b$ of Theorem~\ref{var:rational} through $b=10$.
The verification program checks the complete differential equality
$\Acal_m=P_m+DH_m$ through weight $24$, not only the retained
three-variable sector.

\subsection{Anchored primitives and higher derivative sectors}

For any real $x,x_0>0$, the same identity gives the ordinary anchored
primitive relation
\begin{equation}\label{mom:anchored}
 \int_{x_0}^x P_m(\psi(t),\psi'(t),\ldots)\dd t
 =\int_{x_0}^x\Acal_m(t)\dd t-H_m(\psi(x),\psi'(x),\ldots)
                           +H_m(\psi(x_0),\psi'(x_0),\ldots).
\end{equation}
It states exactly how the primitive changes under normal reduction;
it does not postulate an elementary antiderivative for $\Acal_m$.
Section~\ref{res:section} gives a separate polylogarithmic family
whose repeated Euler primitives and convergent endpoints are explicit.
For additional differential sectors, Lemma~\ref{var:finitejet} applies
without change with $d=3,4,\ldots$. The resulting equations involve
more variables but retain all variational endpoint information.
